\chapter{Die mathematischen Bauteile von TFPT}
\label{ch:algebra}
\section{Die Ausgangsannahmen enthalten mehr als zwei Zahlen}
Die historische Kurzfassung nennt $c_3=1/(8\pi)$ und einen fünfteiligen Träger. Vollständig gelesen setzt P1 einen orientierten Rand beziehungsweise eine Naht, einen geeigneten positiven Kern und eine Normierung voraus. P2 setzt Art und Struktur des Trägers voraus; die spätere $3+2$-Lesart wird mit Farb- und schwachen Richtungen verbunden. Orientierung, Positivität, komplexe Struktur und Markierungen sind echte Daten. Sie verschwinden nicht dadurch, dass eine Zusammenfassung nur zwei Zahlen nennt. \src{S09, S. 4--5; S05, §2}

Der alternative Anker $a=(1,1,2)$ organisiert dieselben kleinen Zahlen. Seine elementarsymmetrischen Größen sind $4,5,2$, seine Potenzsummen
\begin{equation}
p_n=2+2^n,\qquad p_1p_2p_3=4\cdot6\cdot10=240,
\qquad240+p_4-p_3=248.
\end{equation}
Diese Gleichungen sind exakt. Sie sind eine kompakte Buchhaltung, noch keine Konstruktion einer Lie-Algebra: Für $E_8$ braucht man die wirklichen Vektoren, ihre Paarungen und die Operationsstruktur.

\section{Vom Fünferträger zum Spinor}
Bei fünf komplexen Richtungen kann man antisymmetrische Kombinationen von null, zwei und vier Richtungen bilden. Zusammen erhält man
\begin{equation}
S^+=\Lambda^{\mathrm{even}}\C^5,
\qquad\dim S^+=\binom50+\binom52+\binom54=1+10+5=16.
\end{equation}
Das ist die gerade Spinorstruktur von $\mathrm{Spin}(10)$. „Antisymmetrisch“ heißt hier: Beim Vertauschen zweier Plätze wechselt ein Vorzeichen. Es ist eine genaue Bauvorschrift, keine bildliche Ähnlichkeit.

\begin{bild}{Fünf Fächer, mehrere Belegungsmuster}
Man kann sich fünf beschriftete Fächer vorstellen. Man benutzt kein Fach, zwei Fächer oder vier Fächer und zählt die verschiedenen Kombinationen. So entstehen $1+10+5$ Möglichkeiten. Die Vorzeichenregeln zwischen den Fächern tragen die zusätzliche Spinorstruktur. Die Zahl 16 allein wäre dafür nicht ausreichend.
\end{bild}

Der im Korpus verwendete Familienzähler $(16-1)/5=3$ liefert im gewählten Aufbau drei Kopien beziehungsweise 48 Weyl-Komponenten. Er ist für sich kein chiraler Indexsatz. Warum drei propagierende Familien überleben und konjugierte unerwünschte Partner verschwinden, bleibt eine dynamische oder topologische Auswahlfrage.

\section{Warum die Verklebung wirklich zu E\texorpdfstring{$_8$}{8} führt}
Die beiden positiven Wurzelgitter $D_5$ und $A_3\cong D_3$ besitzen Determinante vier und Diskriminantengruppe $\Z_4$. Eine Diskriminantengruppe beschreibt hier die möglichen Restklassen, mit denen man das Gitter ohne Verlust der Ganzzahligkeit ergänzen kann. Für passende Generatoren gilt modulo eins
\begin{equation}
q_{D_5}(k)=\frac{5k^2}{8},\qquad
q_{A_3}(k)=\frac{3k^2}{8}.
\end{equation}
Diagonal gekoppelt ist die Summe $k^2$ ganzzahlig. Die vier gemeinsamen Klassen passen also zu einer geraden Erweiterung. Für den Index vier folgt
\begin{equation}
\det L=\frac{\det D_5\,\det A_3}{4^2}=1.
\end{equation}
Das positive gerade unimodulare Gitter vom Rang acht ist $E_8$. Unter diesen markierten Ausgangsdaten ist die Hülle damit exakt rekonstruiert. \src{S05, §10; S06}

Eine konkrete Realisierung beginnt mit $L_0=D_5\oplus D_3$ und
\begin{equation}
g=(\tfrac12,\ldots,\tfrac12)\in\R^8,
\qquad L=L_0+\Z g.
\end{equation}
Die Klasse von $g$ hat Ordnung vier, $g^2=2$ und ganzzahlige Paarungen mit $L_0$. In der Standardbeschreibung bestehen die Wurzeln aus 112 ganzzahligen Vektoren mit genau zwei Einträgen $\pm1$ und 128 halbzahligen Vektoren $(\pm\frac12)^8$ mit gerader Zahl von Minuszeichen. Alle haben Normquadrat zwei.

\begin{figure}[H]\centering
\begin{tikzpicture}
\node[box,text width=42mm] (d) at (0,0) {$D_5$\\Rang 5 · Determinante 4\\Restklassen $\Z_4$};
\node[box,text width=42mm] (a) at (7,0) {$A_3$\\Rang 3 · Determinante 4\\Restklassen $\Z_4$};
\node[box,text width=90mm] (e) at (3.5,-2.5) {$E_8$\\Passende diagonale Verklebung · Index 4\\Rang 8 · Determinante 1 · 240 Wurzeln};
\draw[arr] (d.south)--(e.north west);
\draw[arr] (a.south)--(e.north east);
\node[font=\sffamily\small,text=slate,align=center] at (3.5,.1) {Quadratische\\Formen ergänzen sich};
\end{tikzpicture}
\caption{Die Aussage $5+3=8$ ist nur die Rangzählung. Erst die passende Verklebung der Restklassen macht aus dem Produkt die gerade unimodulare $E_8$-Hülle.}
\end{figure}

\section{Wurzeln und Generatoren sind verschiedene Zählungen}
Die 240 Wurzeln sind besondere Vektoren in einem acht-dimensionalen Raum. Hinzu kommen acht Cartanrichtungen, sodass die Lie-Algebra 248 Dimensionen besitzt. Unter $D_5\times A_3$ lautet ihre komplexe Zerlegung
\begin{equation}
\mathfrak e_8=(45,1)\oplus(1,15)\oplus(10,6)
\oplus(16,4)\oplus(\overline{16},\overline4).
\label{eq:branch}
\end{equation}
\par\noindent\begin{tabularx}{\linewidth}{L{47mm}rY}
\toprule
\textbf{Sektor}&\textbf{Wurzeln}&\textbf{Generatoren einschließlich Cartan}\\\midrule
$D_5$&40&$45$\\
$A_3$&12&$15$\\
Gemischter Vektorsektor&60&$(10,6)$ mit Dimension $60$\\
Erste Spinorklasse&64&$(16,4)$\\
Konjugierte Spinorklasse&64&$(\overline{16},\overline4)$\\\bottomrule
\end{tabularx}\par

Die konjugierte Spinorklasse muss auf beiden Faktoren korrekt bezeichnet werden. Die Schreibweise $(16,\overline4)$ in Teilen von S08/S04 ist dafür unzureichend, sofern nicht ausdrücklich eine andere Konvention erklärt wird. Die Gesamtdimension allein bemerkt diesen Typfehler nicht.

Für die $\Z_4$-Grade gilt
\begin{align}
\mathfrak g_0&=(45,1)\oplus(1,15),&\mathfrak g_1&=(16,4),\\
\mathfrak g_2&=(10,6),&\mathfrak g_3&=(\overline{16},\overline4),
\end{align}
mit $[\mathfrak g_a,\mathfrak g_b]\subset\mathfrak g_{a+b\bmod4}$. Deshalb führen gleichartige Spinorgrade nach $(10,6)$, konjugierte Grade hingegen nach $\mathfrak g_0$. Die Quellen berichten 960 geordnete Wurzelpaare für $1+1\to2$, sowie $640+192$ Wurzelsummen und 64 Gegenwurzelpaare mit Cartanziel für $1+3\to0$. Diese Typregel ist später für die Kopplung entscheidend.

\section{Das endliche Fenster: 60 Strahlen und 15 Kontexte}
Ein Strahl ist ein reiner Zustand ohne seine insgesamt irrelevante globale Phase. Ein Kontext ist eine vollständige orthogonale Viererbasis: vier mögliche Ausgänge einer scharfen Messung. Im markierten komplexen Viererträger ergeben sich 60 besondere Strahlen in 15 solchen Kontexten. In der gewählten Kodierung sind dies die reinen Zwei-Qubit-Stabilizerzustände und ihre vollständigen Pauli-Kontexte. \src{S09, Kap. 3 und 9}

Die 15 nichttrivialen Pauli-Operatoren sind \emph{nicht} dieselben 15 Objekte wie die Kontexte. Jeder Kontext enthält drei miteinander kommutierende nichttriviale Paulis; jede Pauli-Richtung gehört zu drei Kontexten. Die Geometrie lässt sich mit $\mathbb F_2^4$ und seiner symplektischen Paarung beschreiben; ihre Symmetriegruppe ist $\mathrm{Sp}(4,2)\cong S_6$ mit 720 Elementen.

Die 60 Projektoren spannen den gesamten 16-dimensionalen Raum hermitescher $4\times4$-Operatoren linear auf. Ihre \emph{positiven Mischungen} ergeben jedoch nur das Stabilizer-Polytop. „Jede Richtung linear ausdrücken“ und „jeden Quantenzustand als zulässige Mischung herstellen“ sind verschiedene Aussagen.

\section{Die Räume im Größenvergleich}
\begin{longtable}{L{45mm}L{31mm}L{74mm}}
\toprule
\textbf{Objekt}&\textbf{Größe}&\textbf{Was die Zahl zählt}\\\midrule\endhead
$E_8$-Wurzeln&240&Besondere Vektoren, keine 240 unabhängigen Raumrichtungen.\\
$\mathfrak e_8$&248&Dimension der Generatoralgebra.\\
Ein Viererträger&$\C^4$&Zustände eines Ququarts, hier als zwei Qubits kodiert.\\
Vier solche Träger&$4^4=256$&Komplexe Hilbertraumdimension der Zelle.\\
15 allgemeine CQ-Blöcke&$15\cdot16=240$&Reelle lineare Dichtematrixkoordinaten.\\
16 Majoranas&8 komplexe Moden&Fockraumdimension $2^8=256$.\\
Vier Teilchen in 8 Moden&$\binom84=70$&Ein fester Besetzungssektor des Fockraums.\\
Drei Familien, verdoppelt&96&Eine andere Familien-/Konjugationsbuchhaltung.\\\bottomrule
\end{longtable}
Gleiche Zahlen können verschiedene Objekte beschreiben. Eine Identifikation benötigt eine Abbildung, die Operatoren, Positivität, Phasen, Zusammensetzung und Auslesung erhält.

S11 verschärft die Kritik an der 240-Bijektion: Die natürlichen CQ-Indizes $(C,P)$ zerfallen unter der angegebenen endlichen Wirkung in Bahnen $15+45+180$. Die Wurzel-/Strahlseite besitzt bei der dort vorausgesetzten transitiven Strahlwirkung anders strukturierte Bahnen. Unter genau dieser äquivarianten Zuordnung passt die Bijektion nicht. Das ist stärker als eine bloße Warnung vor Numerologie, aber kein Verbot jeder denkbaren anders definierten Abbildung oder Gruppenwirkung. Die gesamte Orbitprüfung aus S11 wurde hier nicht erneut ausgeführt.
